\subsection{OpenMP -- Open Multi-Processing}
\label{sec:frameworks:openmp}

\begin{quotebox}{\href{https://www.openmp.org/about/openmp-faq/\#WhatIs}{OpenMP FAQ}}
OpenMP is a specification for a set of compiler directives, library routines, and environment variables that can be used to specify high-level parallelism in Fortran and C/\cpp programs.
\end{quotebox}

\acrfull{openmp} is a standard created by the OpenMP Architecture Review Board with the first specification version 1.0 released in October 1998~\cite{openmp_architecture_review_board_openmp_1998}. 
Back then, \gls{openmp} only supported \glspl{cpu}. 
This was changed with version 4.0 released in July 2013~\cite{openmp_architecture_review_board_openmp_2013}, which introduced \gls{openmp} target offloading constructs to also support accelerators like \glspl{gpgpu}. 
The current \gls{openmp} version 6.0 was released in November 2024~\cite{openmp_architecture_review_board_openmp_2024}, showing that \gls{openmp} is still further improved and under active development. 
Since \gls{openmp} is merely a specification, there is no single standard \gls{openmp} implementation but multiple implementations from different vendors targeting different hardware platforms. 
These implementations include but are not limited to implementations from \gls{amd}, Intel, NVIDIA, LLVM, GNU GCC, ARM, etc. 

\gls{openmp} consists of two main components: a runtime library and compiler pragmas or directives. 
The runtime library contains functions to retrieve basic information like the total number of threads or the current thread ID in a parallel region but also to change the behavior of an \gls{openmp} program by, e.g., explicitly setting the number of available threads or specifying the used thread scheduling algorithm. 
The actual parallelization in \gls{openmp} is done through compiler directives. 
These directives have the benefit that they are easy to use and non-intrusive. 
Additionally, since these are only directives if a compiler does not support \gls{openmp}, he is free to ignore them. 
This means that code parallelized with \gls{openmp} can also be compiled with compilers that do not support it, granted that no runtime library functions are used. 
However, the resulting code is then obviously only single-threaded.  
Internally, \gls{openmp} uses a fork-join model where an initial thread spawns several so-called work threads in a parallel region. 
For better performance, in general, \gls{openmp} implementations are based on a custom thread pool. 

\begin{listing}
    \begin{moderncode}{cpp}
#pragma omp parallel for
for (int idx = 0; idx < N; ++idx) {
    Y[idx] = alpha * X[idx] + Y[idx];
}

#pragma omp target map(tofrom: Y[0:N]) map(to: X[0:N])
#pragma omp teams distribute parallel for
for (int idx = 0; idx < N; ++idx) {
    Y[idx] = alpha * X[idx] + Y[idx];
}
    \end{moderncode}
    \caption{%
    A simple DAXPY example using \gls{openmp} (top) and \gls{openmp} target offloading (bottom) (\href{https://godbolt.org/\#g:!((g:!((g:!((h:codeEditor,i:(filename:'1',fontScale:14,fontUsePx:'0',j:2,lang:c\%2B\%2B,selection:(endColumn:32,endLineNumber:25,positionColumn:32,positionLineNumber:25,selectionStartColumn:32,selectionStartLineNumber:25,startColumn:32,startLineNumber:25),source:'\%23include+\%3Ciostream\%3E\%0A\%23include+\%3Cvector\%3E\%0A\%0A\%0Aint+main()+\%7B\%0A++const+int+N+\%3D+1024\%3B\%0A\%0A++//+create+and+fill+used+data\%0A++std::vector\%3Cdouble\%3E+X(N)\%3B\%0A++std::vector\%3Cdouble\%3E+Y(N)\%3B\%0A++for+(int+i+\%3D+0\%3B+i+\%3C+N\%3B+\%2B\%2Bi)+\%7B\%0A++++X\%5Bi\%5D+\%3D+i\%3B\%0A++++Y\%5Bi\%5D+\%3D+2+*+i\%3B\%0A++\%7D\%0A++const+double+alpha+\%3D+0.5\%3B\%0A\%0A++/*\%0A++//+CPU+only\%0A++\%23pragma+omp+parallel+for\%0A++for+(int+idx+\%3D+0\%3B+idx+\%3C+N\%3B+\%2B\%2Bidx)+\%7B\%0A++++++Y\%5Bidx\%5D+\%3D+alpha+*+X\%5Bidx\%5D+\%2B+Y\%5Bidx\%5D\%3B\%0A++\%7D\%0A++*/\%0A\%0A++const+double+*d_X+\%3D+X.data()\%3B\%0A++double+*d_Y+\%3D+Y.data()\%3B\%0A\%0A++//+OpenMP+target+offloading\%0A++\%23pragma+omp+target+map(tofrom:+d_Y\%5B0:N\%5D)+map(to:+d_X\%5B0:N\%5D)\%0A++\%23pragma+omp+teams+distribute+parallel+for\%0A++for+(int+idx+\%3D+0\%3B+idx+\%3C+N\%3B+\%2B\%2Bidx)+\%7B\%0A++++++d_Y\%5Bidx\%5D+\%3D+alpha+*+d_X\%5Bidx\%5D+\%2B+d_Y\%5Bidx\%5D\%3B\%0A++\%7D\%0A\%0A++for+(int+i+\%3D+0\%3B+i+\%3C+10\%3B+\%2B\%2Bi)+\%7B\%0A++++std::cout+\%3C\%3C+Y\%5Bi\%5D+\%3C\%3C+!'+!'\%3B\%0A++\%7D\%0A++return+0\%3B\%0A\%7D'),l:'5',n:'0',o:'C\%2B\%2B+source+\%232',t:'0')),header:(),k:52.66571512767246,l:'4',m:100,n:'0',o:'',s:0,t:'0'),(g:!((g:!((h:compiler,i:(compiler:g142,filters:(b:'0',binary:'1',binaryObject:'1',commentOnly:'0',debugCalls:'1',demangle:'0',directives:'0',execute:'0',intel:'0',libraryCode:'0',trim:'1',verboseDemangling:'0'),flagsViewOpen:'1',fontScale:14,fontUsePx:'0',j:1,lang:c\%2B\%2B,libs:!(),options:'-O3+-std\%3Dc\%2B\%2B17+-fopenmp',overrides:!(),selection:(endColumn:1,endLineNumber:1,positionColumn:1,positionLineNumber:1,selectionStartColumn:1,selectionStartLineNumber:1,startColumn:1,startLineNumber:1),source:2),l:'5',n:'0',o:'+x86-64+gcc+14.2+(Editor+\%232)',t:'0')),k:100,l:'4',m:49.98157361341441,n:'0',o:'',s:0,t:'0'),(g:!((h:output,i:(compilerName:'x86+nvc\%2B\%2B+25.1',editorid:2,fontScale:14,fontUsePx:'0',j:1,wrap:'1'),l:'5',n:'0',o:'Output+of+x86-64+gcc+14.2+(Compiler+\%231)',t:'0')),header:(),l:'4',m:50.01842638658559,n:'0',o:'',s:0,t:'0')),k:47.334284872327544,l:'3',n:'0',o:'',t:'0')),l:'2',n:'0',o:'',t:'0')),version:4}{https://godbolt.org/z/63nT7Eaeo}).
    }
    \label{code:openmp_example}
\end{listing}

\autoref{code:openmp_example} shows a simple DAXPY example in the first four lines. 
Parallelizing a simple for loop with \gls{openmp} is as easy as annotating it with a simple \code{#pragma omp parallel for} directive. 
This is also called loop parallelism in \gls{openmp}. 
Additionally, \gls{openmp} also supports task parallelism, which is not shown here since it will not be used in the remainder of this dissertation. 
\autoref{code:openmp_example} also shows a \gls{openmp} target offloading example of the same code at the bottom. 
While the first code can only be executed on \glspl{cpu}, the target offloading code can also be executed on any supported \gls{gpu}, automatically migrating memory from and to the \gls{gpu} when necessary. 

One downside of \gls{openmp}'s high-level directive-based constructs is that low-level performance optimizations are generally impossible, resulting in worse performance than hand-tuned native codes. 
Since the \gls{openmp} code must run on many different hardware platforms, its high-level \gls{api} can not include low-level optimizations that may be only available on a few hardware platforms. 
In such cases, the programmer has to switch to another framework or hope the compiler translates the high-level construct to the desired low-level instructions. 


%%% ADVANTAGES AND DISADVANTAGES
\AdvantagesAndDisadvantages{openmp}{
    \item standardized
    \item supports a wide variety of hardware platforms
    \item non-intrusive and easy-to-use
}{
    \item very high-level, i.e., some optimizations not possible
}
